%%%%%%%%%%%%%%%%%%%%%%%%%%%%%%%%%%%%%%%%%
% Drake Certificate
% LaTeX Template
% Version 1.0 (February 24, 2022)
%
% This template originates from:
% https://www.LaTeXTemplates.com
%
% Author:
% Vel (vel@latextemplates.com)
%
% License:
% CC BY-NC-SA 4.0 (https://creativecommons.org/licenses/by-nc-sa/4.0/)
%
%%%%%%%%%%%%%%%%%%%%%%%%%%%%%%%%%%%%%%%%%

%----------------------------------------------------------------------------------------
%	PACKAGES AND OTHER DOCUMENT CONFIGURATIONS
%----------------------------------------------------------------------------------------

\documentclass[
	a4paper, % Paper size, use 'a4paper' for A4 (default) or 'letterpaper' for US letter (you may want to adjust margins afterwards)
]{CSDrakeCertificate}

%---------------------------------------------------------------------------------
%	CERTIFICATE INFORMATION
%---------------------------------------------------------------------------------

\heading{Attestazione} % The main certificate heading at the top, typically will be "Certificate", "Award", "Congratulations" or similar. Will be made uppercase automatically.

\subheading{Corso Base di Biostatistica} % The subheading to the main heading, useful for qualifying the main heading, leave empty if not required (i.e. \subheading{})


\recipient{X	Y} % Recipient name, will be made uppercase automatically

\largedescription{ha frequentato con profitto il corso base di Biostatistica presso la Scuola Dottorati dell'Universit\`a \textit{Magna Gr\ae cia} di Catanzaro durante l'anno accademico 2024-2025 (DD 47/2025, Prot. 121/2025) ed ha sostenuto positivamente l'esame finale.} % The large description line under the recipient name, useful for stating what the certificate is for or writing a congratulatory message, leave empty if not required (i.e. \largedescription{})
\smalldescription{Si rilascia come attestazione del conseguimento di 2 Crediti Formativi Universitari.} % The smaller description under the large description, useful for a further classification of the certificate, leave empty if not required (i.e. \smalldescription{})

\leftlinetitle{Germaneto, 4 agosto 2025} % The title under the left line, will be made uppercase automatically and can be left empty if not required (i.e. \leftlinetitle{}) in which case the right line will be horizontally centered (if present)
\leftlineoptional{} % A subtitle to the title under the left line, useful to qualify the title, can be left empty if not required (i.e. \leftlineoptional{})

\rightlinetitle{Massimo Borelli} % The title under the right line, will be made uppercase automatically and can be left empty if not required (i.e. \rightlinetitle{}) in which case the left line will be horizontally centered (if present)
\rightlineoptional{Autenticazione: ZZZ} % A subtitle to the title under the right line, useful to qualify the title, can be left empty if not required (i.e. \rightlineoptional{})

%----------------------------------------------------------------------------------------

\begin{document}

%---------------------------------------------------------------------------------
%	HEADINGS
%---------------------------------------------------------------------------------

\begin{center} % Center horizontally
	{\fontsize{60pt}{60pt}\selectfont\sourceserifproblack{\MakeUppercase{\heading}}\par} % Main heading
	
	\ifdefempty{\subheading}{}{{\fontsize{45pt}{42pt}\selectfont\textbf{\subheading}\par}} % Output subheading if it was specified
\end{center}

\vfill\vfill\vfill % Automatic vertical whitespace for alignment

%---------------------------------------------------------------------------------
%	RECIPIENT
%---------------------------------------------------------------------------------

\begin{tikzpicture}
	\node [
		rectangle, % Node shape
		fill=BlueGreen, % Background color
		text=white, % Text color
		align=center, % Text alignment
		text width=0.85\textwidth, % Maximum text width, shorter than the full line width so there is horizontal padding
		font=\fontsize{46pt}{46pt}\selectfont, % Font styling
		inner xsep=1.5cm, % Horizontal padding for the triangles
		inner ysep=0.5cm, % Vertical padding
	] (namebox) at (0,0){%
		{\hyphenpenalty 10000\exhyphenpenalty 10000\sourcesanspro{\textbf{\MakeUppercase{\recipient}}}\par}% Output recipient name (suppressing hyphenation)
	};
	\draw [fill=white, draw=white] (namebox.north west) -- (namebox.south west) -- ($(namebox.west)+(0.75, 0)$) -- cycle; % Left triangle
	\draw [fill=white, draw=white] (namebox.north east) -- (namebox.south east) -- ($(namebox.east)+(-0.75, 0)$) -- cycle; % Right triangle
\end{tikzpicture}

\vfill % Automatic vertical whitespace for alignment

%---------------------------------------------------------------------------------
%	DESCRIPTIONS
%---------------------------------------------------------------------------------

\begin{center}	
	\ifdefempty{\largedescription}{}{{\LARGE\largedescription\vspace*{0.02\textheight}}} % Output large description if it was entered
	
	\ifdefempty{\smalldescription}{}{{\large\smalldescription}} % Output small description if it was entered
\end{center}

\vfill\vfill\vfill\vfill\vfill\vfill\vfill % Automatic vertical whitespace for alignment

%---------------------------------------------------------------------------------
%	SIGNATURE/INFORMATION LINES
%---------------------------------------------------------------------------------

\begin{center}	
	\ifdefempty{\leftlinetitle}{}{ % Display the left line and title(s) if the left title has been entered
		\begin{tikzpicture}
			\node [rectangle, fill=black, minimum width=0.35\textwidth, inner ysep=2pt] (leftline) at (0, 0){}; % Left line
			\node [rectangle, anchor=north, minimum height=20pt] at ($(leftline.south)-(0, 3pt)$) {{\LARGE\sourcesanspro{\textbf{\MakeUppercase{\leftlinetitle}}}}}; % Top title
			\node [rectangle, anchor=north, minimum height=20pt] at ($(leftline.south)-(0, 20pt)$) {{\large\sourcesanspro{\textit{\leftlineoptional}}}}; % Subtitle
		\end{tikzpicture}%
	}
	\ifdefempty{\leftlinetitle}{}{\ifdefempty{\rightlinetitle}{}{\hfill}} % If both lines have titles and are therefore visible, space them out horizontally
	\ifdefempty{\rightlinetitle}{}{ % Display the right line and title(s) if the right title has been entered
		\begin{tikzpicture}
			\node [rectangle, fill=black, minimum width=0.35\textwidth, inner ysep=2pt] (rightline) at (0, 0){}; % Right line
			\node [rectangle, anchor=north, minimum height=20pt] at ($(rightline.south)-(0, 3pt)$) {{\LARGE\sourcesanspro{\textbf{\MakeUppercase{\rightlinetitle}}}}}; % Top title
			\node [rectangle, anchor=north, minimum height=20pt] at ($(rightline.south)-(0, 20pt)$) {{\large\sourcesanspro{\textit{\rightlineoptional}}}}; % Subtitle
		\end{tikzpicture}%
	}
\end{center}

%---------------------------------------------------------------------------------
%	CERTIFICATE FRAMING
%---------------------------------------------------------------------------------

\outputframe % Output the certificate frames and corner accents

%---------------------------------------------------------------------------------

\end{document}
